\chapter{Hardware and Software Environment}\label{hardware-and-software-environment}

The entire pipeline, including the fine-tuning of the Tier-1 classifier, was developed and run on a single consumer laptop. The constrained resources, in particular the 4 GB of GPU memory, directly shaped two design choices discussed in Chapter~\ref{methods-and-implemented-modules}: the selection of a distilled transformer for Tier-1, and the training batch size of 8. Core count and memory matter too: throughput rose as classifier processes were added up to the four physical cores, and memory sat close to full once the container stack shared the machine (Section~\ref{throughput-scalability}). Table~\ref{tab:env-hardware} gives the full specification.

\section{Hardware}\label{appendix-hardware}

\begin{longtable}{@{}
  >{\raggedright\arraybackslash}p{\dimexpr 0.4000\linewidth - \tabcolsep \relax}
  >{\raggedright\arraybackslash}p{\dimexpr 0.6000\linewidth - \tabcolsep \relax}@{}}
\caption{Hardware environment.}\label{tab:env-hardware}\tabularnewline
\toprule
Component & Specification \\
\midrule
\endfirsthead
\toprule
Component & Specification \\
\midrule
\endhead
\bottomrule
\endlastfoot
Processor (CPU) & Intel Core i5-8265U, 8th generation, 4 physical cores and 8 logical threads \\
System memory & 12 GB DDR4 RAM (8 GB + 4 GB, 2400 MHz) \\
Graphics (GPU) & NVIDIA GTX 1650 with Max-Q Design, 4 GB VRAM \\
Operating system & Windows 11 Pro \\
Tier-1 inference & CPU by default; the GPU is selectable through one configuration value, and both were measured (Section~\ref{throughput-scalability}) \\
Tier-2 inference & \texttt{gpt-oss:120b} (Ollama tag \texttt{gpt-oss:120b-cloud}) served through Ollama's hosted endpoint, a 116.8-billion-parameter model far larger than the local GPU can hold \\
\end{longtable}

\section{Software}\label{appendix-software}

The pipeline runs on \textbf{Python 3.10.11}. Table~\ref{tab:env-python} lists the Python packages that affect a reported result, and Table~\ref{tab:env-services} the containerised services, each pinned to the exact image tag used.

\begin{longtable}{@{}
  >{\raggedright\arraybackslash}p{\dimexpr 0.3000\linewidth - 1.33\tabcolsep \relax}
  >{\raggedright\arraybackslash}p{\dimexpr 0.2000\linewidth - 1.33\tabcolsep \relax}
  >{\raggedright\arraybackslash}p{\dimexpr 0.5000\linewidth - 1.33\tabcolsep \relax}@{}}
\caption{Python runtime packages.}\label{tab:env-python}\tabularnewline
\toprule
Package & Version & Role \\
\midrule
\endfirsthead
\toprule
Package & Version & Role \\
\midrule
\endhead
\bottomrule
\endlastfoot
PyTorch (\texttt{torch}) & 2.5.1 (CUDA 12.1) & Deep-learning framework for DistilBERT \\
Transformers & 5.1.0 & DistilBERT model and tokenizer \\
scikit-learn & 1.7.2 & Logistic Regression and SVM baselines, TF-IDF \\
sentence-transformers & 5.5.1 & MiniLM embedder for RAG and the LLM-Wiki \\
ollama (client) & 0.6.1 & Tier-2 LLM access \\
qdrant-client & 1.17.0 & Vector store client \\
kafka-python & 2.3.0 & Kafka producer and consumer \\
elasticsearch (client) & 7.17.9 & Elasticsearch access \\
streamlit & 1.58.0 & Operator dashboard and analyst chat \\
pandas / numpy / scipy & 2.3.3 / 2.2.6 / 1.15.3 & Data handling, metrics, McNemar's test \\
plotly & 6.8.0 & Dashboard charts \\
\end{longtable}

\begin{longtable}{@{}
  >{\raggedright\arraybackslash}p{\dimexpr 0.4000\linewidth - \tabcolsep \relax}
  >{\raggedright\arraybackslash}p{\dimexpr 0.6000\linewidth - \tabcolsep \relax}@{}}
\caption{Containerised services, as pinned in \texttt{docker-compose.yml}.}\label{tab:env-services}\tabularnewline
\toprule
Service & Image tag \\
\midrule
\endfirsthead
\toprule
Service & Image tag \\
\midrule
\endhead
\bottomrule
\endlastfoot
Apache Kafka (Confluent) & \texttt{confluentinc/cp-kafka:7.5.0} \\
Apache Zookeeper & \texttt{zookeeper:3.9} \\
Apache Spark (master and worker) & \texttt{apache/spark:3.4.1} \\
Elasticsearch (server) & \texttt{docker.elastic.co/\allowbreak elasticsearch/\allowbreak elasticsearch:8.10.2} \\
Kibana & \texttt{docker.elastic.co/kibana/kibana:8.10.2} \\
Qdrant & \texttt{qdrant/qdrant:v1.10.1} \\
\end{longtable}

One pairing is deliberate and worth flagging for anyone rebuilding the stack: the Elasticsearch \emph{client} is a 7.17 release while the \emph{server} is 8.10, which the client supports in its compatibility mode. Upgrading the client to match the server major version is the obvious-looking correction, and it breaks the ingestion path.

\section{Models}\label{appendix-models}

Table~\ref{tab:env-models} lists the three models, with their size where it constrained a design choice.

\begin{longtable}{@{}
  >{\raggedright\arraybackslash}p{\dimexpr 0.3400\linewidth - 1.33\tabcolsep \relax}
  >{\raggedright\arraybackslash}p{\dimexpr 0.3000\linewidth - 1.33\tabcolsep \relax}
  >{\raggedright\arraybackslash}p{\dimexpr 0.3600\linewidth - 1.33\tabcolsep \relax}@{}}
\caption{Models used, with size where it constrained a design choice.}\label{tab:env-models}\tabularnewline
\toprule
Model & Size & Role \\
\midrule
\endfirsthead
\toprule
Model & Size & Role \\
\midrule
\endhead
\bottomrule
\endlastfoot
\texttt{distilbert-base-uncased}, fine-tuned & 67.0M parameters, 6 layers, about 269 MB on disk & Tier-1 classifier \\
\texttt{gpt-oss:120b} via Ollama & 116.8B parameters (5.1B active per token, mixture-of-experts), hosted & Tier-2 judge and the four agents \\
\texttt{all-MiniLM-L6-v2} & 384-dimensional embeddings & Retrieval for both RAG and the LLM-Wiki \\
\end{longtable}
